% Card: rear checks (rook endings)
\paragraph{Rear checks}\index{Rook endings!Rear checks}

\diagram[rear-checks]{4k3/8/4P3/4K3/8/8/6R1/r7 b - - 1 1}

Here the pawn has already reached the sixth rank, so the black rook can no longer
\term{cut-off} the white king on that rank, as in \xref{diagram}{cut-off}. Instead it goes to the
far end of the board and checks the king from behind. These are
\vocab[rook endings]{rear-checks}{rear checks}{The defending rook checks the attacking king from
behind, from the far end of the board, along the files. The king cannot hide in front of its pawn,
because the defending king stands there}:\index{Practical play!Clear-cut draw}

\moves{1...\R e1+}

Wherever the white king goes, the rook checks it again from behind on its new file:
\moves{2.\K d6 \R d1+} or \moves{2.\K f6 \R f1+}. The king cannot hide in front of its pawn,
because the black king blocks the way, and as soon as it walks away from the pawn, the black king
takes it. This is a clear-cut draw.

As with \term{side-checks}, the rook must have enough room to give checks: there must be three
ranks between the pawn and the rook. That is why the defending rook often goes to the first or
second rank.
